% Appendices A--C.
\chapter{构建与运行}\label{app:build}

本节给出从干净仓库复现实验结果的最小命令集。所有命令在仓库根目录执行，构建目录为 \code{build/}。

\begin{lstlisting}[language=bash,basicstyle=\ttfamily\small]
# 配置、编译与测试
cmake -S . -B build
cmake --build build -j
ctest --test-dir build --output-on-failure

# Task 1：词频统计
./build/mse count --csv build/wordcount.csv --top 50

# Bonus：熵驱动停用词
./build/mse stopwords --derive --min-df-ratio 0.4 --min-entropy-rate 0.93 \
                      --compare data/stopwords.txt

# Task 2：正式索引（v1）
./build/mse build --index build/index.bin \
                  --stop-words data/stopwords.txt --stem

# Bonus：v2 压缩索引 + 磁盘 B+Tree 词典
./build/mse build --index build/index-v2.bin \
                  --stop-words data/stopwords.txt --stem --compress
./build/mse build --index build/index.bin \
                  --stop-words data/stopwords.txt --stem \
                  --dict-file build/index.bpt

# Bonus：SPIMI 块式构建
./build/mse build-blocks --stop-words data/stopwords.txt --stem \
                         --block-terms 4096

# Bonus：B+Tree 规模实测
./build/mse bptree-bench --keys 10000000 --lookups 1000

# Task 3：查询
./build/mse query --index build/index.bin \
                  --stop-words data/stopwords.txt --word osric --pretty
./build/mse query --index build/index.bin \
                  --stop-words data/stopwords.txt --phrase "my lord"
./build/mse query --index build/index.bin \
                  --stop-words data/stopwords.txt --words "king queen" --min-match 2

# Bonus：磁盘词典查询与全词表校验
./build/mse dump --index build/index.bin --dict-file build/index.bpt --term banquo
./build/mse dump --index build/index.bin --dict-file build/index.bpt --verify-all

# Bonus：第二语料 RFC
./build/mse stats --manifest data2/data/manifest.tsv
./build/mse count --manifest data2/data/manifest.tsv --csv build/rfc_wordcount.csv
./build/mse stopwords --derive --manifest data2/data/manifest.tsv \
                      --min-df-ratio 0.4 --min-entropy-rate 0.93 \
                      --out build/rfc_stopwords.txt
./build/mse build --manifest data2/data/manifest.tsv \
                  --stop-words build/rfc_stopwords.txt --stem \
                  --index build/rfc_index.bin

# Task 4：阈值实验
python3 tools/threshold_matrix.py

# 独立对拍
python3 tools/check_wordcount.py > build/wordcount_py.csv
diff build/wordcount.csv build/wordcount_py.csv
python3 tools/brute_force.py --stop-words data/stopwords.txt --stem
python3 tools/check_phrases.py --index build/index.bin \
    --stop-words data/stopwords.txt --stem --queries 300

# Bonus：500k 文件合成跑批（仅用于规模实验，不参与正式语料）
python3 tools/gen_synthetic_corpus.py --root build/scale-500k \
    --docs 500000 --tokens-per-doc 40 --vocab 5000000
./build/mse build-blocks --manifest build/scale-500k/data/manifest.tsv \
    --block-terms 50000 --index build/scale-500k-index.bin
\end{lstlisting}

若系统没有 \code{cmake}，可用 g++ 直接编译；仓库的 CMake 配置也支持开启“警告即错误”（选项名 \code{MSE_WERROR}）。Python 脚本只依赖标准库（Porter 词干化使用 vendored 的 NLTK 扩展模式数据，见 \code{tools/}）。

\chapter{数据与索引格式速查}\label{app:formats}

\section{manifest}

\code{data/manifest.tsv} 为制表符分隔、首行表头，共 8 列（第 8 列可选）：

\begin{lstlisting}[language=,basicstyle=\ttfamily\small]
docid  kind  category  title  text  source  url  [format]
30     play  Tragedy   Hamlet data/text/0030_hamlet.txt  data/raw/hamlet/full.html  http://shakespeare.mit.edu/hamlet/full.html
\end{lstlisting}

\code{format} 取 \code{marked}（默认，带行标记）或 \code{plain}（每行都是正文）。\code{data2/} 的 RFC 语料使用 \code{plain}，因此其中的方括号和圆括号不会被当作舞台提示删除。

\section{文本层行标记}

\begin{table}[htbp]
  \centering
  \caption{\code{data/text/*.txt} 的行标记}
  \label{tab:line-markers}
  \small
  \begin{tabularx}{0.8\linewidth}{@{}l X l@{}}
    \toprule
    标记 & 含义 & 默认 \\
    \midrule
    \code{[ACT I]} / \code{[SCENE ...]} & 标题 & 过滤 \\
    \code{(Exit Ghost)} & 舞台指示 & 过滤 \\
    \code{@BERNARDO} & 说话人 & 过滤 \\
    其他 & 正文 & 保留 \\
    \code{#} & 文件头注释 & 跳过 \\
    \bottomrule
  \end{tabularx}
\end{table}

\section{倒排索引文件}

索引文件 v1/v2 的布局见第 2.6.4 节：magic 为 \code{MSE1}，小端；version 1 存储原始 u32 postings，version 2 存储差分加 varint postings；头部、文档表、term 表布局相同。磁盘 B+Tree 词典 sidecar 的 magic 为 \code{BPT1}，4 KiB 页，页 0 为头部，其后是树页（见第 2.6.6 节）。

\chapter{评分对照}\label{app:grading}

\begin{table}[htbp]
  \centering
  \caption{题目要求与本报告的对应位置}
  \label{tab:grading}
  \small
  \begin{tabularx}{\linewidth}{@{}X r X@{}}
    \toprule
    评分项 & 分值 & 本报告对应位置 \\
    \midrule
    Word counting 程序 & 1 & 2.4、3.5；代码 \code{src/count.cpp} \\
    Index generation 程序 & 5 & 2.6、3.6；代码 \code{src/index.cpp} 等（含 \code{bptree}、\code{blockbuild}） \\
    Query processing 程序 & 3 & 2.7、3.7；代码 \code{src/query.cpp} \\
    倒排索引正确性测试 & 2 & 3.2--3.4、3.9 \\
    查询阈值测试 & 2 & 3.8 \\
    分析与评论 & 3 & 第 4 章 \\
    Bonus & +2 & 4.7（熵停用词、RFC、v2 压缩、B+Tree、SPIMI、500k 文件、50M 键） \\
    文档第 1 章 & 1 & 第 1 章 \\
    文档第 2 章 & 2 & 第 2 章 \\
    整体文档风格 & 1 & 全文结构、表格、术语统一 \\
    \bottomrule
  \end{tabularx}
\end{table}
